\subsection{FAMILIES MULTIPLIABLE IN R*}

In the multiplicative group \(\mathbb{R}^*\) of finite non-zero real numbers, a family \((x_{t})_{i\in I}\) can be multipliable only if \(\lim x_{t} = 1\) with respect to the filter of complements of finite subsets of I (Chapter III, § 5, no. 2, Proposition 1). In particular there can be only a finite number of indices \(i\) such that \(x_{t} < 0\) . We may therefore limit ourselves to considering only families \(x_{t}\) all of whose terms are strictly positive; it is then convenient to put \(x_{i} = i + u_{i}\) , where the \(u_{i}\) are subject to the inequalities \(-i < u_{i} < +\infty\) for all \(i\) . Since each point of \(\mathbf{R}^*\) has a countable fundamental system of neighbourhoods, the set of indices \(i\) such that \(u_{i} \neq 0\) is countable, if the family \((i + u_{i})\) is multipliable in \(\mathbf{R}^*\) .

THEOREM 4. The family \((1 + u_{t})\) is multipliable in \(\mathbb{R}^*\) if and only if the family \((u_{t})\) is summable in \(\mathbb{R}\) .

Lemma. (i) If \((a_{i})_{1\leqslant i\leqslant p}\) is a finite sequence of numbers \(>0\) , then

\[
\prod_ {i = 1} ^ {p} \left(1 + a _ {i}\right) \geqslant 1 + \sum_ {i = 1} ^ {p} a _ {i}.\tag{2}
\]

\begin{enumerate}
\def\labelenumi{(\roman{enumi})}
\setcounter{enumi}{1}
\tightlist
\item
  If also \(a_i < 1\) for each \(i\) , then
\end{enumerate}

\[
\prod_ {i = 1} ^ {p} \left(\mathrm{I} - a _ {i}\right) \geqslant \mathrm{I} - \sum_ {i = 1} ^ {p} a _ {i}.\tag{3}
\]

These inequalities are clear if \(p = 1\) , and are proved by induction on \(p\) . If

\[
\prod_ {i = 1} ^ {p - 1} \left(\mathrm{I} + a _ {i}\right) \geqslant \mathrm{I} + \sum_ {i = 1} ^ {p - 1} a _ {i},
\]

we have

\[
\prod_ {i = 1} ^ {p} \left(\mathrm{I} + a _ {i}\right) \geqslant \left(\mathrm{I} + a _ {p}\right) \left(\mathrm{I} + \sum_ {i = 1} ^ {p - 1} a _ {i}\right) = \mathrm{I} + \sum_ {i = 1} ^ {p} a _ {i} + a _ {p}. \sum_ {i = 1} ^ {p - 1} a _ {i} \geqslant \mathrm{I} + \sum_ {i = 1} ^ {p} a _ {i}.
\]

Again, if

\[
\prod_ {i = 1} ^ {p - 1} \left(\mathrm{I} - a _ {i}\right) \geqslant \mathrm{I} - \sum_ {i = 1} ^ {p - 1} a _ {i},
\]

we have

\[
\prod_ {i = 1} ^ {p} \left(\mathrm{I} - a _ {i}\right) \geqslant \left(\mathrm{I} - a _ {p}\right) \left(\mathrm{I} - \sum_ {i = 1} ^ {p - 1} a _ {i}\right) = \mathrm{I} - \sum_ {i = 1} ^ {p} a _ {i} + a _ {p}. \sum_ {i = 1} ^ {p - 1} a _ {i} \geqslant \mathrm{I} - \sum_ {i = 1} ^ {p} a _ {i}.
\]

Having proved the lemma, we note that if the family \((\mathbf{I} + u_{t})\) is multipliable then so are the families \((\mathbf{I} + u_{t}^{+})\) and \((\mathbf{I} - u_{t}^{-})\) , since \(\mathbf{R}^*\) is a complete group (Chapter III, § 5, no. 3, Proposition 2); and conversely, if the families \((\mathbf{I} + u_{t}^{+})\) and \((\mathbf{I} - u_{t}^{-})\) are multipliable, then so is \((\mathbf{I} + u_{t})\) (Chapter III, § 5, no. 3, Proposition 3). We need therefore consider only the cases in which all the \(u_{t}\) are \(\geqslant 0\) and in which they are all \(\leqslant 0\) .

Suppose first that \(u_{t} \geqslant 0\) for all \(t\) . If the family \((1 + u_{t})\) is multipliable, then for each \(\varepsilon > 0\) there is a finite subset \(J\) of the index set in \(I\) such that, for each finite subset \(H\) of \(I\) disjoint from \(J\) , we have \(\mathbf{I} \leqslant \prod_{i \in \mathbf{H}} (\mathbf{I} + u_i) \leqslant \mathbf{I} + \varepsilon\) ; by (2) it follows that \(\sum_{i \in \mathbf{H}} u_i \leqslant \varepsilon\) , which shows that \((u_i)\) is summable in \(\mathbf{R}\) by virtue of Cauchy's criterion (Chapter III, § 5, no. 2, Theorem I).

Conversely, suppose that \((u_{t})\) is summable in \(\mathbf{R}\) . For each \(\varepsilon\) such that \(0 < \varepsilon < 1\) , there is a finite subset \(J\) of \(I\) such that, for each finite subset \(H\) of \(I\) disjoint from \(J\) , we have \(0 \leqslant \sum_{t \in H} u_{t} \leqslant \varepsilon\) . By (3), we have therefore \(\prod_{t \in H} (I - u_{t}) \geqslant I - \varepsilon\) ; but since \(I + u \leqslant \frac{I}{I - u}\) for every number \(u\) such that \(0 \leqslant u < 1\) , it follows that

\[
\mathrm{I} \leqslant \prod_ {i \in \mathrm{H}} (\mathrm{I} + u _ {i}) \leqslant \frac {\mathrm{I}}{\mathrm{I} - \varepsilon},
\]

and this shows that \((\mathbf{I} + u_{t})\) is multipliable (Cauchy's criterion).

The proof is similar when all the \(u_{i}\) are \(\leqslant 0\) . To show that \((u_{i})\) is summable if \((1 + u_{i})\) is multipliable, use formula (2) and the inequality \(1 - u \leqslant \frac{1}{1 + u}\) ( \(0 \leqslant u < 1\) ); to show that \((1 + u_{i})\) is multipliable if \((u_{i})\) is summable, use formula (3).

In Chapter V (§ 4), the topological study of the group \(\mathbf{R}^*\) will enable us to give another criterion for multipliability of a family in \(\mathbf{R}^*\) with the help of the logarithmic function; in a later volume we shall establish the equivalence of this criterion with the one above, using the differential properties of the logarithm.

